\begingroup\fontsize{9.3}{10.7}\selectfont\setlength{\parskip}{1pt}
\section{Notación de las realizaciones gravitatorias}
La notación distingue los objetos de la construcción, los operadores
de respuesta y sus realizaciones físicas. Los productos internos positivos
de las extensiones y la métrica lorentziana tienen funciones diferentes.
\renewcommand{\arraystretch}{1.0}\begin{longtable}{@{}p{.23\linewidth}p{.72\linewidth}@{}}
\toprule
Símbolo & Objeto y función\\ \midrule
\endhead
\(G_a,\hbar_a\) & Secciones recíprocas de gravitación y acción.
Su producto determina el área invariante
\(\ell_P^2=\hbar_aG_a/c^3=L_\alpha^2\) en la realización radial.\\
\(L_*,L_\alpha=L\) & Sección longitudinal de referencia y longitud
construida \(L_\alpha=(5759/23040)\alpha^{16}L_*\).\\
\(E_L,t_P,t_0\) & Energía de retorno \(E_L=\hbar_{\rm ret}c/L_\alpha\),
tiempo \(t_P=L_\alpha/c\) y paso \(t_0=\pi_{\rm HMT}t_P/54\).\\
\(R_X,\Lambda_X\) & Lector radial positivo y longitud conjugada del
mismo carácter; \(R_X\Lambda_X=L_\alpha^2I\).\\
\(R_h,\zeta_h\) & Radio areal de horizonte y coordenada
\(R_h/L_\alpha=\zeta_h>0\). Conservan la condición geométrica del horizonte.\\
\(\chi_{\rm rad}=L_\alpha/E_L\) & Proporción dimensional entre radio
y energía. Sus unidades y su función difieren de las del registro \(K\)
y de \(\kappa=8\pi G/c^4\).\\
\(\gamma_{\rm HMT}\) & Evaluación del funcional hexada--octada,
reutilizada como coeficiente del operador de Holst.\\
\(k_B,\Theta\) & Coeficiente de conversión informacional y temperatura
en la sección térmica declarada.\\
\(P_5,\Gamma_5,\Lambda_n\) & Proyección sobre las cinco coordenadas,
isometría tensorial y reducción transversal de rango dos.\\
\(c_s,c_g\) & Cociclo de memoria traslacional y componente de retorno
del transporte orientado.\\
\(\beta_m^\square,\Lambda_m^\square\) & Soldadura de pantalla y
respuesta energética, con su identidad cuadrática y su escala.\\
\(\mathsf C,\mathsf L,f_{\min}\) & Incidencia transportada,
Gramiano y extensión de mínima energía para el dato de frontera.\\
\(s,\sigma,\mathsf M_{\rm cel}\) & Covector del diferencial, corriente
material y emparejamiento que convierte sus coordenadas.\\
\(e,\omega,\mathring\omega,K,T\) & Cotetrada, conexión, conexión de
Levi--Civita, contorsión y torsión.\\
\(\mathsf P_\gamma,\mathsf L_e\) & Operador interno de Holst y
aplicación lineal de torsión a la ecuación de conexión.\\
\(Q,F,\Delta S\) & Término gravitatorio cuadrático, acción material
y corrección efectiva tras eliminar la contorsión.\\
\(\mathsf H\) & Hessiano de la acción compuesta que controla la
continuación local de una rama torsional.\\
\(M_n,a,a_{\rm ref}\) & Memoria, factor de escala y normalización
de referencia de su realización homogénea.\\
\(s_0,\rho_0,w\) & Amplitud torsional, densidad material de
referencia y parámetro barotrópico del dominio de rebote.\\
\(\varrho,V_c,\widehat{\mathscr Q}_{\rm sp}\) & Estado material normalizado, volumen
comóvil y observable cuártico normalmente ordenado; determinan
el funcional \(s_0[\varrho,V_c]\).\\
\(C_\gamma\) & Coeficiente de la interacción efectiva espinorial
obtenido al eliminar la contorsión.\\
\(\tau,t_{\rm b}\) & Duración característica e instante del mínimo
de la solución homogénea de polvo.\\
\(F_0,R,a_{\rm b}\) & Función de Friedmann, perturbación y
umbral transversal de la evolución homogénea.\\
\(T_\Lambda,T_0\) & Partes de traza y sin traza del tensor residual;
sus divergencias registran el intercambio entre sectores.\\
\(\mathsf{G}_{\rm ret},\mathsf{G}_{\rm adv}\) & Propagadores retardado y avanzado.
Aquí la letra \(G\) designa operadores, no la constante gravitatoria.\\
\(P_{\rm abs}\) & Proyección sobre las fuentes que satisfacen la
igualdad de las dos lecturas causales en la frontera.\\
\bottomrule
\end{longtable}

\par\endgroup
